\subsection{Rayleigh Quotients and Minimax}

Throughout this subsection, $n\in\mathbb{Z}_{>0}$.

\begin{definition} \label{def:increasing-hermitian-eigenvalue-ordering} \hlt{Hermitian Eigenvalue Order and Rayleigh Quotient}\\
Let $A\in\operatorname{Mat}(n,n,\mathbb{C})$ be Hermitian. Matrix $A$ has real eigenvalues and admits an orthonormal eigenbasis by Theorems \ref{thm:self-adjoint-eigenvalues-real} and \ref{thm:complex-spectral-theorem}.
\begin{enumerate}[label=(\roman*)]
\setlength{\itemsep}{0pt}
\item Write eigenvalues with algebraic multiplicity as $\lambda_1(A)\leq\lambda_2(A)\leq\cdots\leq\lambda_n(A)$. If $A$ is fixed, abbreviate $\lambda_i(A)$ to $\lambda_i$.
\item For $\mathbf{x}\in\mathbb{C}^n\setminus\{\mathbf{0}\}$, the \hlt{Rayleigh quotient} is $\operatorname{Ray}_A(\mathbf{x})\defeq\frac{\mathbf{x}^*A\mathbf{x}}{\mathbf{x}^*\mathbf{x}}$. If $\norm{\mathbf{x}}_2=1$, then $\operatorname{Ray}_A(\mathbf{x})=\mathbf{x}^*A\mathbf{x}$.
\end{enumerate}
\end{definition}

\begin{theorem} \label{thm:rayleigh-quotient-spectral-interval} \hlt{Rayleigh Quotient on a Spectral Interval}\\
Let $A\in\operatorname{Mat}(n,n,\mathbb{C})$ be Hermitian with increasing eigenvalues $\lambda_1\leq\cdots\leq\lambda_n$ and corresponding orthonormal eigenvectors $\mathbf{u}_1,\ldots,\mathbf{u}_n$.\\
If $p,q\in J_n$, $p\leq q$, and $\mathbf{x}\in\operatorname{span}(\mathbf{u}_p,\ldots,\mathbf{u}_q)$ satisfies $\norm{\mathbf{x}}_2=1$, then $\lambda_p\leq\operatorname{Ray}_A(\mathbf{x})\leq\lambda_q$.
\begin{proof}
Write $\mathbf{x}=\sum\limits_{i=p}^{q}c_i\mathbf{u}_i$; then $(\forall\,i\in\{p,\ldots,q\})\,\abs{c_i}^2\geq0$ and $\sum\limits_{i=p}^{q}\abs{c_i}^2=1$ by orthonormality.\\
The identity $\operatorname{Ray}_A(\mathbf{x})=\sum\limits_{i=p}^{q}\lambda_i\abs{c_i}^2$ holds by the eigenvector relations and Definition \ref{def:increasing-hermitian-eigenvalue-ordering}(ii).\\
The bound $\lambda_p\leq\operatorname{Ray}_A(\mathbf{x})\leq\lambda_q$ holds by $\lambda_p\leq\lambda_i\leq\lambda_q$ and the preceding weight relations.\\
\hlt{Structural Commentary}: Eigenbasis expansion reduces the quotient to a weighted average over its spectral support.\\
\hlt{Key Insight}: The spectral weights $\abs{c_i}^2$ are nonnegative and sum to one.
\end{proof}
\end{theorem}

\begin{corollary} \label{cor:rayleigh-extrema-and-diagonal-entry-bounds} \hlt{Rayleigh Extrema and Diagonal Entry Bounds}\\
Let $A=[a_{ij}]\in\operatorname{Mat}(n,n,\mathbb{C})$ be Hermitian with orthonormal eigenpairs $(\lambda_i,\mathbf{u}_i)_{i\in J_n}$, where $\lambda_1\leq\cdots\leq\lambda_n$.\\
For $p,q\in J_n$ with $p\leq q$, set $\mathcal{U}_{p,q}\defeq\operatorname{span}(\mathbf{u}_p,\ldots,\mathbf{u}_q)$.
\begin{enumerate}[label=(\roman*)]
\setlength{\itemsep}{0pt}
\item $\min\limits_{\substack{\mathbf{x}\in\mathcal{U}_{p,q}\\ \norm{\mathbf{x}}_2=1}}\operatorname{Ray}_A(\mathbf{x})=\lambda_p$.
\item $\max\limits_{\substack{\mathbf{x}\in\mathcal{U}_{p,q}\\ \norm{\mathbf{x}}_2=1}}\operatorname{Ray}_A(\mathbf{x})=\lambda_q$.
\item $(\forall\,\mathbf{x}\in\mathbb{C}^n)\,[\norm{\mathbf{x}}_2=1\Longrightarrow\lambda_1\leq\operatorname{Ray}_A(\mathbf{x})\leq\lambda_n]$.
\item $(\forall\,i\in J_n)\,\lambda_1\leq a_{ii}\leq\lambda_n$.
\end{enumerate}
\begin{proof}
\begin{enumerate}[label=(\roman*)]
\setlength{\itemsep}{0pt}
\item The lower bound and its attainment at $\mathbf{u}_p$ hold by Theorem \ref{thm:rayleigh-quotient-spectral-interval} and $A\mathbf{u}_p=\lambda_p\mathbf{u}_p$, respectively.
\item The upper bound and its attainment at $\mathbf{u}_q$ hold by Theorem \ref{thm:rayleigh-quotient-spectral-interval} and $A\mathbf{u}_q=\lambda_q\mathbf{u}_q$, respectively.
\item The asserted bound holds by (i)--(ii) with $(p,q)=(1,n)$ and $\mathcal{U}_{1,n}=\mathbb{C}^n$.
\item $(\forall\,i\in J_n)\,[\norm{\mathbf{e}_i}_2=1\land\operatorname{Ray}_A(\mathbf{e}_i)=a_{ii}]$ by Proposition \ref{prop:coordinate-norms-on-fn}(ii), Definition \ref{def:increasing-hermitian-eigenvalue-ordering}(ii), and coordinate multiplication.\\
The relation $(\forall\,i\in J_n)\,\lambda_1\leq a_{ii}\leq\lambda_n$ holds by the preceding identity and (iii).
\end{enumerate}
\hlt{Structural Commentary}: Endpoint eigenvectors attain support bounds; coordinate vectors bound every diagonal entry.\\
\hlt{Key Insight}: Attainment turns spectral enclosures into exact extrema.
\end{proof}
\end{corollary}

Figure \ref{fig:rayleigh-quotient-convex-combination} places the quotient inside the eigenvalue interval determined by its spectral support, as asserted by Theorem \ref{thm:rayleigh-quotient-spectral-interval}.

\begin{illustrationfigure}[H]
  \centering
  \begin{tikzpicture}[
  >=Stealth,
  scale=0.88,
  axis/.style={->, gray!65, line width=0.65pt},
  eig/.style={circle, fill=blue!65!black, inner sep=1.7pt},
  rq/.style={circle, fill=red!75!black, inner sep=2.2pt},
  lab/.style={font=\small, align=center, fill=white, inner sep=1pt},
  tiny/.style={font=\small, align=center}
]
  \draw[axis] (-0.4,0) -- (7.4,0);
  \foreach \x/\labx in {0.8/{\lambda_1},2.2/{\lambda_p},5.0/{\lambda_q},6.4/{\lambda_n}} {
    \draw[gray!55] (\x,0.09) -- (\x,-0.09);
    \node[eig] at (\x,0) {};
    \node[tiny, below=4pt] at (\x,0) {$\labx$};
  }
  \draw[green!55!black, line width=1.25pt] (2.2,0.55) -- (5.0,0.55);
  \node[lab, above=4pt] at (3.6,0.55) {$\mathcal{U}_{p,q}$};
  \node[rq] at (3.7,0) {};
  \node[lab, below=16pt] at (3.7,0) {$\operatorname{Ray}_A(\mathbf{x})$};
\end{tikzpicture}

  \caption{For ordered orthonormal eigenpairs of a Hermitian $A$, a unit vector supported on $\mathbf{u}_p,\ldots,\mathbf{u}_q$ has $\operatorname{Ray}_A(\mathbf{x})\in[\lambda_p,\lambda_q]$.}
  \label{fig:rayleigh-quotient-convex-combination}
\end{illustrationfigure}

\begin{computationmethod} \label{meth:rayleigh-eigenbasis-method} \hlt{Rayleigh Eigenbasis Method}\\
Given a Hermitian $A\in\operatorname{Mat}(n,n,\mathbb{C})$ and a unit vector $\mathbf{x}\in\mathbb{C}^n$, construct a certified spectral-support enclosure for $\operatorname{Ray}_A(\mathbf{x})$ as follows:
\begin{enumerate}[label=(\roman*)]
\setlength{\itemsep}{0pt}
\item compute an orthonormal eigenbasis $\mathbf{u}_1,\ldots,\mathbf{u}_n$ in increasing eigenvalue order by Theorem \ref{thm:complex-spectral-theorem} and Definition \ref{def:increasing-hermitian-eigenvalue-ordering}(i);
\item compute $c_i\defeq\mathbf{u}_i^*\mathbf{x}$ for $i\in J_n$ and the nonempty support $I(\mathbf{x})\defeq\{i\in J_n:c_i\neq0\}$;
\item evaluate $\operatorname{Ray}_A(\mathbf{x})=\sum\limits_{i\in J_n}\lambda_i\abs{c_i}^2$;
\item set $p\defeq\min I(\mathbf{x})$ and $q\defeq\max I(\mathbf{x})$, and output the certified enclosure $\lambda_p\leq\operatorname{Ray}_A(\mathbf{x})\leq\lambda_q$ by Theorem \ref{thm:rayleigh-quotient-spectral-interval}.
\end{enumerate}
\end{computationmethod}

\begin{theorem} \label{thm:courant-fischer-rayleigh-minmax} \hlt{Rayleigh Minimax Principle}\\
Let $A=A^*\in\operatorname{Mat}(n,n,\mathbb{C})$ have eigenvalues $\lambda_1\leq\cdots\leq\lambda_n$. For $k\in J_n$:
\begin{enumerate}[label=(\roman*)]
\setlength{\itemsep}{0pt}
\item $\lambda_k=\min\limits_{\substack{\mathcal{S}\leq\mathbb{C}^n\\ \dim(\mathcal{S})=k}}\ \max\limits_{\mathbf{0}\neq\mathbf{x}\in\mathcal{S}}\operatorname{Ray}_A(\mathbf{x})$.
\item $\lambda_k=\max\limits_{\substack{\mathcal{T}\leq\mathbb{C}^n\\ \dim(\mathcal{T})=n-k+1}}\ \min\limits_{\mathbf{0}\neq\mathbf{x}\in\mathcal{T}}\operatorname{Ray}_A(\mathbf{x})$.
\end{enumerate}
\begin{proof}
Fix an orthonormal eigenbasis $\mathbf{u}_1,\ldots,\mathbf{u}_n$ in increasing eigenvalue order by Theorem \ref{thm:complex-spectral-theorem} and Definition \ref{def:increasing-hermitian-eigenvalue-ordering}(i).\\
Set $E_k\defeq\operatorname{span}(\mathbf{u}_1,\ldots,\mathbf{u}_k)$ and $F_k\defeq\operatorname{span}(\mathbf{u}_k,\ldots,\mathbf{u}_n)$. $(\forall\,\mathbf{x}\in\mathbb{C}^n\setminus\{\mathbf{0}\})\,\operatorname{Ray}_A(\mathbf{x})=\operatorname{Ray}_A(\mathbf{x}/\norm{\mathbf{x}}_2)$ by Definition \ref{def:increasing-hermitian-eigenvalue-ordering}(ii).\\
Hence $\max\limits_{\mathbf{0}\neq\mathbf{x}\in E_k}\operatorname{Ray}_A(\mathbf{x})=\lambda_k$ and $\min\limits_{\mathbf{0}\neq\mathbf{x}\in F_k}\operatorname{Ray}_A(\mathbf{x})=\lambda_k$ by Corollary \ref{cor:rayleigh-extrema-and-diagonal-entry-bounds}(ii),(i).
\begin{enumerate}[label=(\roman*)]
\setlength{\itemsep}{0pt}
\item If $\mathcal{S}\leq\mathbb{C}^n$ and $\dim(\mathcal{S})=k$, then $\mathcal{S}\cap F_k\neq\{\mathbf{0}\}$ by Corollary \ref{cor:large-subspaces-intersect}(i) and $\dim(\mathcal{S})+\dim(F_k)=n+1$.\\
The bound $\max\limits_{\mathbf{0}\neq\mathbf{x}\in\mathcal{S}}\operatorname{Ray}_A(\mathbf{x})\geq\lambda_k$ holds by the preceding intersection, Theorem \ref{thm:rayleigh-quotient-spectral-interval}, and scale invariance.\\
The outer minimum equals $\lambda_k$ by this bound and $\mathcal{S}=E_k$ in the preceding extremum identity.
\item If $\mathcal{T}\leq\mathbb{C}^n$ and $\dim(\mathcal{T})=n-k+1$, then $\mathcal{T}\cap E_k\neq\{\mathbf{0}\}$ by Corollary \ref{cor:large-subspaces-intersect}(i) and $\dim(\mathcal{T})+\dim(E_k)=n+1$.\\
The bound $\min\limits_{\mathbf{0}\neq\mathbf{x}\in\mathcal{T}}\operatorname{Ray}_A(\mathbf{x})\leq\lambda_k$ holds by the preceding intersection, Theorem \ref{thm:rayleigh-quotient-spectral-interval}, and scale invariance.\\
The outer maximum equals $\lambda_k$ by this bound and $\mathcal{T}=F_k$ in the preceding extremum identity.
\end{enumerate}
\hlt{Structural Commentary}: Spectral-flag intersections supply universal bounds; $E_k$ and $F_k$ attain them.\\
\hlt{Key Insight}: A dimension sum of $n+1$ forces a shared comparison vector.
\end{proof}
\end{theorem}
